\documentclass[10pt,a4paper,sans]{moderncv}
\moderncvstyle{classic}
\moderncvcolor{blue}
\usepackage[utf8]{inputenc}
\usepackage[scale=0.85]{geometry}
\usepackage{helvet}

\name{Alexandro}{Disla}
\title{MEAL Department Head --- Monitoring, Evaluation, Accountability and Learning}
\address{Rue Saint-Louis Jeanty \#14, Route de Frère}{Port-au-Prince, Haiti}{}
\phone[mobile]{(+509) 4148-3700}
\email{alexandrodisla@hotmail.com}
\extrainfo{GitHub: \url{https://github.com/AD0791}}

\begin{document}
\makecvtitle

\section{Professional Profile}
\cvitem{}{More than ten years across monitoring, evaluation and information systems --- \textbf{five years of NGO MEAL work} in three Haitian sectors (health and HIV, water and sanitation, education) and \textbf{eight years of public investment monitoring and statistical modelling} at the Economic and Social Planning Directorate of Haiti's Ministry of Planning. My work is about what makes a programme's numbers defensible: logical frameworks and indicator tracking tables, data quality assurance protocols established and then supervised, digital collection standardised so that data is correct at entry rather than repaired afterwards. An applied economist by training and a practising data engineer, I build and administer the databases, pipelines and dashboards a country office depends on --- I do not merely specify them. Based in Port-au-Prince, available for field travel; fluent in French, English and Haitian Creole.}

\section{Core Competencies}
\cvitem{}{
\begin{itemize}
    \item \textbf{MEAL framework leadership \& harmonisation:} M\&E strategies and frameworks, results chains, logical frameworks, indicator tracking tables (ITT), MEAL plans, baselines, targets and required disaggregations --- including sex and age disaggregation decided at indicator definition. Standardisation of forms and data flows across sites and sectors.
    \item \textbf{Data quality \& control:} DQA protocols established and supervised, verification back to source registers, corrective actions tracked to closure; beneficiary list deduplication and reliable counts for donor reporting and proposals.
    \item \textbf{Information systems \& donor reporting:} Relational databases and ETL pipelines designed, administered and secured (role-based access control, secure cloud storage); weekly monitoring dashboards; production and submission of PEPFAR/MER indicators to USAID through the DHIS2/DATIM platforms \textit{(as a data producer)}; project under United Nations (UNOPS) contract.
    \item \textbf{Supervision, learning \& coordination:} Supervision of enumerators on multi-site programmes, technical line management of M\&E staff, design and delivery of training on collection tools and data quality; coordination with government partners (MSPP, DINEPA, Ministry of Finance) and translation of technical findings for non-technical leadership.
\end{itemize}}

\section{MEAL and Programme Quality Experience}

\cventry{Jan 2026 -- Mar 2026}{Data Analyst \& MEAL Consultant}{Anseye Pou Ayiti}{Haiti}{}{
\begin{itemize}
    \item Defined the collection strategy and implemented performance indicators for three educational leadership programmes, from indicator definition through to reporting back.
    \item Coded XLSForm instruments for ODK and Google Forms with skip logic and validation constraints, so that errors are stopped at entry rather than corrected downstream.
    \item Administered the ODK system and admin panel, maintained Python integration scripts on AWS, and trained field agents on the collection tools.
\end{itemize}}

\cventry{Oct 2023 -- Jan 2026}{mWater Consultant \& Data Analyst}{HANWASH}{Haiti / Remote}{}{
\begin{itemize}
    \item \textbf{MEAL design:} Refined MEAL plans and indicator tracking tables aligned to JMP (WHO/UNICEF) standards for water, hygiene and sanitation programmes.
    \item \textbf{Data management:} Deployed complex geospatial surveys on mWater, administered the database and resolved field collection incidents to preserve the integrity of the series.
    \item \textbf{Automation \& capacity:} Built automated dashboards (mWater, Power BI) connected to collection APIs for real-time reporting; wrote the technical guides and trained teams to use them independently.
\end{itemize}}

\cventry{Oct 2021 -- Aug 2024}{Monitoring \& Evaluation Officer}{Impact Youth Project, Caris Foundation International}{Haiti}{}{
\begin{itemize}
    \item \textbf{Collection systems:} Established digital collection (CommCare, HIVHaiti) on a multi-site health and HIV programme, standardising forms and data flows across sites.
    \item \textbf{Quality assurance (DQA):} Established and supervised Data Quality Assessment protocols on integrated MySQL data --- tracing reported figures back to source registers and following corrective actions through to closure rather than to mere logging.
    \item \textbf{Donor reporting:} Produced and submitted the PEPFAR/MER indicators destined for USAID through DHIS2/DATIM, with the consistency checks and reporting cycle that this imposes.
    \item \textbf{Supervision:} Supervised enumerators across all sites and provided technical line management to M\&E staff on collection procedures and quality controls.
    \item \textbf{Automation:} Developed activity reporting in Python and SQL, cutting manual processing time by 40 \%; maintained programme dashboards on AWS (Quarto) and Power BI for weekly indicator monitoring.
\end{itemize}}

\cventry{Nov 2015 -- Aug 2023}{Economist --- Economic and Social Planning Directorate (DPES)}{Ministry of Planning and External Cooperation (MPCE)}{Port-au-Prince}{}{
\begin{itemize}
    \item Contributed to the execution monitoring reports of the 2014--2016 Three-Year Investment Plan: indicator monitoring of public investment projects across every sector and geographic department.
    \item Assigned to the statistics and modelling unit, and contributed to building the ``Development Planning Model'' --- a forecasting, simulation and public policy impact analysis tool.
    \item Took part in macroeconomic framing meetings with the Ministry of Economy and Finance and monitored national macroeconomic indicators.
\end{itemize}}

\section{Data Engineering and Systems Experience}

\cventry{Feb 2026 -- Present}{Software Engineer (Fullstack)}{Tekkod LLC}{Remote}{}{Led a mobile infrastructure modernisation programme with offline reliability and data security for field use; built Looker Studio dashboards with advanced search and filtering for operational visibility.}

\cventry{Mar 2021 -- May 2024}{Backend Developer \& Data Engineer}{Tekkod LLC}{Remote}{}{Built ETL pipelines (Apache Beam) migrating MongoDB data to BigQuery and administered relational databases under concurrent access; delivered secured REST services (FastAPI, JWT) with role-based access control and full OpenAPI documentation.}

\cventry{Nov 2015 -- Mar 2021}{Software Consultant \& Data Integration Specialist}{CassionSoft (UNOPS Project)}{Haiti}{}{Data integration work on a project under United Nations contract, with direct exposure to UN procedures and documentation standards; implemented RBAC authentication (OAuth2/JWT) and validated integration logic through unit testing.}

\section{Education}
\cvitem{2010--2014}{Graduate studies in Applied Economics --- C.T.P.E.A (Centre de Techniques de Planification et d'Économie Appliquée), Port-au-Prince. \textit{Coursework completed; the Diplôme d'Études Supérieures awaits the defence of the exit thesis --- official attestation enclosed with this application.}}
\cvitem{2009--2010}{Baccalaureate (with distinction) --- Institution Saint-Louis de Gonzague, Port-au-Prince}
\cvitem{Self-study}{FlatIron School (Software Engineering), CodeWithMosh, Udemy}

\section{Technical Skills}
\cvitem{MEAL tools}{\textbf{KoboToolbox, ODK, XLSForm}, CommCare, mWater, HIVHaiti, DHIS2/DATIM \textit{(data submission)}, logical frameworks, ITT, DQA protocols}
\cvitem{Analysis}{\textbf{Advanced Excel}, \textbf{Power BI (expert)}, SQL, Python (Pandas), R (RMarkdown/Quarto), Looker Studio, statistical modelling, sampling and trend analysis}
\cvitem{Infrastructure}{MySQL, PostgreSQL, MongoDB, Azure, AWS, Google Cloud (BigQuery), Docker, Git}
\cvitem{Languages}{\textbf{French (native)}, \textbf{English (fluent)}, \textbf{Haitian Creole (native)}}

\end{document}
